% --- GLOSSAREINTRÄGE ---

\newglossaryentry{Multi-Faktor-Authentisierung}{
    name={Multi-Faktor-Authentisierung},
    description={Ein Sicherheitsverfahren, bei dem mehrere unabhängige Faktoren zur Authentifizierung kombiniert werden.}
}

\newglossaryentry{passwort}{
    name={Passwort},
    description={Eine geheime Zeichenkette, die zur Authentisierung dient.}
}

\newglossaryentry{entropie}{
    name={Entropie},
    description={Ein Mass aus der Informationstheorie für die Stärke und Zufälligkeit eines Passworts.}
}

\newglossaryentry{bruteforce}{
    name={Brute-Force-Angriff},
    description={Ein automatisierter Angriff, bei dem Programme alle möglichen Zeichenkombinationen systematisch \mbox{durchprobieren}.}
}

\newglossaryentry{dictionaryattack}{
    name={Dictionary Attack (Wörterbuchangriff)},
    description={Ein gezielter Angriff, bei dem Listen mit \mbox{gebräuchlichen} Wörtern, Namen und bekannten Passwörtern systematisch getestet werden.}
}

\newglossaryentry{datenleck}{
    name={Datenleck},
    description={Ein Vorfall, bei dem vertrauliche Daten unabsichtlich oder absichtlich veröffentlicht oder gestohlen werden.}
}

\newglossaryentry{passwortmanager}{
    name={Passwort-Manager},
    description={Software zur sicheren Verwaltung und Generierung von Passwörtern.}
}

\newglossaryentry{passphrase}{
    name={Passphrase},
    description={Eine längere Zeichenkette, oft bestehend aus mehreren zufälligen Wörtern, die als Passwort verwendet wird.}
}

\newglossaryentry{phishing}{
    name={Phishing},
    description={Eine Form des Social Engineering, bei der Angreifer versuchen, über gefälschte E-Mails, Nachrichten oder Webseiten an vertrauliche Daten zu gelangen.}
}

\newglossaryentry{socialengineering}{
    name={Social Engineering},
    description={Die Manipulation von Personen durch psychologische Tricks, um an vertrauliche Informationen zu gelangen.}
}

\newglossaryentry{spearphishing}{
    name={Spear-Phishing},
    description={Ein gezielter Phishing-Angriff, der auf eine spezifische Person oder bestimmte Mitarbeiter zugeschnitten ist.}
}

\newglossaryentry{quishing}{
    name={Quishing},
    description={Eine Phishing-Methode, bei der manipulierte QR-Codes eingesetzt werden, um Opfer auf gefälschte Webseiten zu leiten.}
}
